% ============================================================
% f2014_sichuan_science_q08_cube.tex
% 2014 四川卷（理科）第 8 题：正方体中的线面角
% 来源 PDF：普通高考/2014/2014四川理.pdf
% 原页：第 1 页（A3 横排）中栏上部，印面 919，第 8 题题干上方
% 说明：正方体 ABCD-A1B1C1D1，O 为 BD 中点（实心点）；图中画出
%       截面三角形 A1BD（A1B 为实线，A1D、BD 为虚线）。
% ============================================================
\begin{tikzpicture}[>=Stealth, line width=0.7pt, font=\small,
  x=1cm, y=1cm, line cap=round, line join=round]

  \coordinate (A)  at (0,0);
  \coordinate (B)  at (2,0);
  \coordinate (C)  at (2.86,0.60);
  \coordinate (D)  at (0.86,0.60);
  \coordinate (A1) at (0,2);
  \coordinate (B1) at (2,2);
  \coordinate (C1) at (2.86,2.60);
  \coordinate (D1) at (0.86,2.60);
  \coordinate (O)  at ($(B)!0.5!(D)$);

  % 可见棱
  \draw (A) -- (B) -- (C) -- (C1) -- (D1) -- (A1) -- cycle;
  \draw (A) -- (A1);
  \draw (B) -- (B1);
  \draw (B1) -- (C1);
  \draw (A1) -- (B1);

  % 不可见棱
  \draw[dashed] (A) -- (D);
  \draw[dashed] (D) -- (C);
  \draw[dashed] (D) -- (D1);

  % 截面三角形 A1BD
  \draw (A1) -- (B);
  \draw[dashed] (A1) -- (D);
  \draw[dashed] (B) -- (D);

  % 点 O
  \fill (O) circle (1.6pt);
  \node at ($(O)+(-0.24,-0.15)$) {$O$};

  % 顶点标注
  \node[below left] at (A) {$A$};
  \node[below right] at (B) {$B$};
  \node[right] at (C) {$C$};
  \node[left=1pt, yshift=2pt] at (D) {$D$};
  \node[left] at (A1) {$A_1$};
  \node[above=1pt] at (B1) {$B_1$};
  \node[right] at (C1) {$C_1$};
  \node[above left] at (D1) {$D_1$};
\end{tikzpicture}
